% ============================================================================
% Plantilla de informe -- Laboratorio de Circuitos Eléctricos II (LIE-311)
% La portada se agrega aparte; este archivo comienza después de ella.
%
% Reglas de entrega del sílabo: PDF con todos los apartados; en prácticas
% virtuales o híbridas, adjuntar también los montajes .ms14 dentro del .zip o
% .rar. El nombre solicitado es LNIE221_NOMBRE_APELLIDO_NUMEROCUENTA.pdf.
% Reemplace todos los marcadores \rellenar{...} antes de entregar.
% ============================================================================

\documentclass[12pt,a4paper]{article}
% ============================================================================
% Formato.tex: configuración común de los informes del Laboratorio de
% Circuitos Eléctricos II (LIE-311). Se incluye desde Documento.tex.
% La portada se omite intencionalmente y se agrega por separado.
% ============================================================================

% Codificación, idioma y tipografía. inputenc mantiene compatibilidad con
% pdfLaTeX; en LuaLaTeX/XeLaTeX puede omitirse.
\usepackage[utf8]{inputenc}
\usepackage[T1]{fontenc}
\usepackage[spanish,es-nodecimaldot]{babel}
\usepackage{lmodern}

% Matemáticas, unidades, tablas e imágenes.
\usepackage{amsmath,amssymb}
\usepackage{siunitx}                 % Ej.: \SI{10}{\volt}
\usepackage{graphicx}
\usepackage{booktabs}
\usepackage{array}
\usepackage{enumitem}
\usepackage{float}
\usepackage{caption}
\usepackage{xcolor}

% Diseño de página, separación de párrafos y enlaces.
\usepackage{geometry}
\usepackage{fancyhdr}
\usepackage{parskip}
\usepackage{csquotes}
\usepackage[hidelinks]{hyperref}

% El estilo numérico produce citas compactas, apropiadas para IEEE. Para APA,
% cambie style=numeric-comp por style=apa (y agregue el paquete correspondiente
% si su instalación de biblatex no lo incluye).
\usepackage[
    backend=biber,
    style=numeric-comp,
    sorting=none
]{biblatex}

% La comprobación permite compilar la plantilla antes de crear references.bib.
% Al entregar, ese archivo debe contener el manual y todas las fuentes usadas.
\IfFileExists{references.bib}{\addbibresource{references.bib}}{}

\geometry{a4paper, margin=2.5cm, headheight=15pt}
\setlength{\parindent}{0pt}
\setlength{\parskip}{6pt}
\setlist{itemsep=2pt, topsep=3pt}
\captionsetup{font=small, labelfont=bf}
\sisetup{locale=DE, per-mode=symbol, detect-all}

% Datos editables. Se redefinen en Documento.tex con \renewcommand. Se declaran
% antes del encabezado porque este último los utiliza.
\newcommand{\numeroPractica}{N}
\newcommand{\nombrepractica}{Título de la práctica}

% Encabezado y pie. Los valores de la práctica se definen en Documento.tex.
\pagestyle{fancy}
\fancyhf{}
\lhead{LIE-311 -- Laboratorio de Circuitos Eléctricos II}
\rhead{Práctica \numeroPractica}
\cfoot{\thepage}

% Inserta una figura centrada y numerada.
% Uso: \imagen[etiqueta]{archivo}{pie de figura}
% Las capturas virtuales deben mostrar fecha y hora; el archivo puede ser PNG,
% JPG o PDF. Una etiqueta permite referenciar la figura con \ref{fig:montaje}.
\newcommand{\imagen}[3][]{%
    \begin{figure}[H]
        \centering
        \includegraphics[width=0.88\textwidth]{#2}
        \caption{#3}
        \if\relax\detokenize{#1}\relax\else\label{#1}\fi
    \end{figure}%
}

% Entorno reutilizable para objetivos y conclusiones numerados.
\newenvironment{itemsinforme}
    {\begin{enumerate}[label=\arabic*., leftmargin=*]}
    {\end{enumerate}}

% Marcador visible para localizar contenido pendiente. Debe reemplazarse o
% eliminarse antes de entregar el informe final.
\newcommand{\rellenar}[1]{\textcolor{red}{\textbf{[Completar: #1]}}}

\hypersetup{
    pdftitle={Informe de laboratorio -- \nombrepractica},
    pdfauthor={Nombre del estudiante},
    pdfsubject={Laboratorio de Circuitos Eléctricos II -- LIE-311}
}


% Cambie estos dos valores para cada práctica.
\renewcommand{\numeroPractica}{1}
\renewcommand{\nombrepractica}{Osciloscopio, generador de funciones y voltaje senoidal}

\begin{document}

% La portada debe contener: institución/logo, título de la práctica, clase
% teórica y sección, catedrático, sección del laboratorio, instructor,
% estudiante, número de cuenta y fecha de entrega.

\section{Introducción}
% Escriba una introducción propia (no la del manual) de 120--150 palabras.
% Presente el tema, su importancia y su conexión con el marco teórico; revise
% redacción y ortografía.
\rellenar{introducción de 120--150 palabras}

\section{Objetivos}
% Incluya como mínimo tres objetivos. Se pueden usar los del manual, pero deben
% corresponder a lo que efectivamente se realizó.
\begin{itemsinforme}
    \item \rellenar{objetivo específico 1}
    \item \rellenar{objetivo específico 2}
    \item \rellenar{objetivo específico 3}
\end{itemsinforme}

\section{Materiales y equipo}
% Liste todos los materiales, instrumentos, componentes, software y equipo
% utilizados. Añada lo que no aparezca en el manual; incluya modelo, rango o
% valor nominal cuando ayude a repetir la práctica.
\begin{itemize}
    \item \rellenar{material, instrumento o componente}
    \item \rellenar{material, instrumento o componente}
    \item \rellenar{software o simulador, si aplica}
\end{itemize}

\section{Marco teórico}
% Desarrolle el marco teórico completo indicado en el manual: conceptos, leyes,
% ecuaciones y condiciones de operación. Explique variables y unidades; no lo
% reduzca a una lista de fórmulas. Cite las fuentes utilizadas.
\subsection{Conceptos fundamentales}
\rellenar{teoría relevante de la práctica, con citas}

\subsection{Ecuaciones y relaciones utilizadas}
\begin{equation}
    \text{\rellenar{ecuación principal}}
    \label{eq:principal}
\end{equation}
\rellenar{definición de variables, unidades y condiciones de aplicación}

\section{Procedimiento experimental}
% Describa en orden cada paso, conexión, configuración y condición de medición.
% Incluya capturas claras de montajes, mediciones y procedimientos. En trabajo
% virtual, las capturas deben mostrar fecha y hora.
\subsection{Montaje y configuración}
\begin{enumerate}
    \item \rellenar{paso de montaje o configuración 1}
    \item \rellenar{paso de montaje o configuración 2}
    \item \rellenar{paso de montaje o configuración 3}
\end{enumerate}

% Ejemplo: \imagen[fig:montaje]{imagenes/montaje.png}{Montaje experimental.}

\subsection{Mediciones}
% Use unidades en los encabezados y documente cualquier condición especial.
\begin{table}[H]
    \centering
    \caption{Registro de mediciones experimentales.}
    \label{tab:mediciones}
    \begin{tabular}{@{}c c c c@{}}
        \toprule
        Ensayo & Magnitud 1 & Magnitud 2 & Observaciones \\
        \midrule
        1 & \rellenar{valor} & \rellenar{valor} & \rellenar{nota} \\
        2 & \rellenar{valor} & \rellenar{valor} & \rellenar{nota} \\
        \bottomrule
    \end{tabular}
\end{table}

\section{Análisis de resultados}
% Es la sección con mayor ponderación (50%). Desarrolle todos los apartados y
% preguntas del manual; muestre cálculos completos, sustitución de valores y
% unidades. Compare teoría, simulación y medición, explique diferencias y cite
% las investigaciones solicitadas.
\subsection{Cálculos y desarrollo matemático}
\begin{align}
    \text{\rellenar{cálculo completo, con unidades}}
    \label{eq:calculo}
\end{align}

\subsection{Comparación y discusión}
\rellenar{comparación de resultados teóricos, simulados y experimentales}

\subsection{Respuestas a las preguntas del manual}
\begin{enumerate}
    \item \textbf{Pregunta 1.} \rellenar{respuesta justificada}
    \item \textbf{Pregunta 2.} \rellenar{respuesta justificada}
\end{enumerate}

\section{Conclusiones}
% Escriba como mínimo tres conclusiones amplias y coherentes. Relaciónelas con
% los objetivos y la evidencia experimental; no repita simplemente el método.
\begin{itemsinforme}
    \item \rellenar{conclusión relacionada con el objetivo 1 y los resultados}
    \item \rellenar{conclusión relacionada con el objetivo 2 y los resultados}
    \item \rellenar{conclusión relacionada con el objetivo 3 y los resultados}
\end{itemsinforme}

\section{Bibliografía}
% Registre el manual y todas las fuentes en APA o IEEE. Cite dentro del texto
% con \cite{clave}; las claves deben existir en references.bib.
\IfFileExists{references.bib}{
    \printbibliography[heading=none]
}{
    \textit{Agregue references.bib y sus fuentes antes de entregar.}
}

\end{document}
